\documentclass[../main]{subfiles}
\begin{document}
\chapter{Idea de Proyecto}
\img{images/idea.jpg}{La idea de proyecto es el motor que nos
  acompa\u00f1a desde el inicio del proceso creativo. Conocemos el origen
  \u2014la idea\u2014 y visualizamos el final: la materializaci\u00f3n f\u00edsica y
funcional de nuestra propuesta.}
\info{Contenido}{
  \begin{itemize}
    \item Definici\u00f3n y alcance de la idea de proyecto.
    \item Metodolog\u00eda de diagn\u00f3stico t\u00e9cnico.
    \item Dise\u00f1o de la soluci\u00f3n base optimizada.
    \item Estructuraci\u00f3n de la idea mediante el Project Idea Canvas.
  \end{itemize}
}
\info{Soluci\u00f3n Base Optimizada}{
  Una soluci\u00f3n base optimizada es aquella dise\u00f1ada espec\u00edficamente
  para satisfacer los requisitos del proyecto con la m\u00e1xima
  eficiencia. Esto requiere un an\u00e1lisis exhaustivo de requerimientos,
  restricciones t\u00e9cnicas y tecnolog\u00edas disponibles.
  Para optimizar una propuesta en el campo electromec\u00e1nico, consideramos:
  \begin{itemize}
    \item \textbf{Rendimiento y Escalabilidad:} Maximizar el sistema
      pensando en futuros crecimientos.
    \item \textbf{Seguridad y Mantenimiento:} Proteger
      usuarios/equipos y facilitar actualizaciones para minimizar
      tiempos de inactividad.
    \item \textbf{Costo y Sostenibilidad:} Minimizar el costo total
      de propiedad a largo plazo con una soluci\u00f3n econ\u00f3mica y duradera.
  \end{itemize}
}
\section{Conceptualizaci\u00f3n de la Idea}
La idea de proyecto es la propuesta inicial para resolver un problema
t\u00e9cnico o aprovechar una oportunidad de mejora. No surge del azar, sino de un proceso de observaci\u00f3n y an\u00e1lisis del entorno industrial o social.

\subsection{Metodolog\u00eda de Diagn\u00f3stico: El \u00c1rbol de Problemas}
Para que una idea sea viable, primero debemos entender la ra\u00edz del problema. El \textbf{\u00c1rbol de Problemas} es una herramienta visual que nos permite organizar la informaci\u00f3n de la siguiente manera:
\begin{itemize}
  \item \textbf{Tronco}: El problema central (lo que observamos).
  \item \textbf{Ra\u00edces}: Las causas que originan el problema (por qu\u00e9 sucede).
  \item \textbf{Ramas/Hojas}: Los efectos o consecuencias del problema (qu\u00e9 pasa si no se resuelve).
\end{itemize}

\subsection{Estructuraci\u00f3n: El Project Idea Canvas}
Una vez diagnosticado el problema, la idea debe formalizarse. En lugar de una simple descripci\u00f3n, utilizamos un \textbf{Canvas de Idea de Proyecto} que debe cubrir:
\begin{description}
  \item[Diagn\u00f3stico:] Resumen del problema basado en el \u00c1rbol de problemas.
  \item[Metas T\u00e9cnicas:] Objetivos cuantificables (ej. "reducir el tiempo de ciclo en 20\%").
  \item[Alcance y L\u00fmites:] Qu\u00e9 har\u00e1 la m\u00e1quina y qu\u00e9 no har\u00e1.
  \item[Restricciones:] Normativas de seguridad, presupuesto m\u00e1ximo y espacio disponible.
  \item[Impacto:] Beneficio esperado para el usuario o la industria.
\end{description}

\newpage
\actividad{Taller de Ideaci\u00f3n de Proyectos}

\textbf{Objetivo}: Transformar una observaci\u00f3n en una propuesta t\u00e9cnica viable.

\begin{enumerate}
  \item \textbf{Fase de Observaci\u00f3n}: Identifique un problema t\u00e9cnico real en el taller de la escuela, en una industria local o en su comunidad.
  \item \textbf{Fase de Diagn\u00f3stico}: Elabore un \textbf{\u00c1rbol de Problemas} donde identifique la causa ra\u00edz y los efectos negativos de dicha situaci\u00f3n.
  \item \textbf{Fase de Propuesta}: Complete la estructura del \textbf{Project Idea Canvas} definiendo las metas t\u00e9cnicas y las restricciones del proyecto.
  \item \textbf{Validaci\u00f3n (The Pitch)}: Presente su Canvas al profesor en una breve exposici\u00f3n (3 minutos) para validar la viabilidad t\u00e9cnica y obtener el consentimiento para iniciar el diseño.
\end{enumerate}

\end{document}